% =====================================================================
%  CHAPITRE 4 - TRAVAUX RÉALISÉS
%  Méthodologie de travail + CS BAG v2 + Elium Sauvegarde + Access TTC
% =====================================================================
\chapter{Travaux réalisés}
\label{chap:travaux}

Ce chapitre présente d'abord la méthodologie de travail appliquée au sein du
service, puis détaille les trois projets sur lesquels j'ai travaillé~: la
refonte de \csbag, l'outil de sauvegarde Elium et le système de rapports
d'accès TTC/GO.

% -----------------------------------------------------------------
\section{Méthodologie de travail au sein de VESI}
\label{sec:methodologie}

Chaque évolution applicative suit un processus structuré en plusieurs étapes,
depuis l'expression du besoin jusqu'à la mise en production et le suivi.
La figure~\ref{fig:workflow} synthétise ce cycle.

\begin{figure}[H]
  \centering
  \scalebox{0.82}{%
  \begin{tikzpicture}[
    font=\scriptsize,
    step/.style={rounded corners=2pt, draw=vinciblue, line width=0.6pt,
      fill=vinciblue!6, align=center, text width=3.2cm, minimum height=0.55cm,
      inner sep=2pt},
    arr/.style={-{Stealth[length=1.8mm]}, line width=0.7pt, vinciblue!80}
  ]
    % Colonne gauche (étapes 1-5)
    \node[step] (s1) {\textbf{1.} Spécifications\\(analyste fonctionnel)};
    \node[step, below=0.2cm of s1] (s2) {\textbf{2.} Réunion de cadrage\\(AF + chef de projet)};
    \node[step, below=0.2cm of s2] (s3) {\textbf{3.} Développement\\(environnement local)};
    \node[step, below=0.2cm of s3] (s4) {\textbf{4.} Déploiement DEV\\(ressources Azure)};
    \node[step, below=0.2cm of s4] (s5) {\textbf{5.} Tests développeur\\(environnement DEV)};
    % Colonne droite (étapes 6-10)
    \node[step, right=1cm of s1] (s10) {\textbf{10.} Monitoring};
    \node[step, below=0.2cm of s10] (s9) {\textbf{9.} Mise en production};
    \node[step, below=0.2cm of s9] (s8) {\textbf{8.} Tests finaux (AF)\\+ validation};
    \node[step, below=0.2cm of s8] (s7) {\textbf{7.} Corrections\\des retours};
    \node[step, below=0.2cm of s7] (s6) {\textbf{6.} Mise en recette\\+ tests AF};
    % Flèches colonne gauche
    \draw[arr] (s1) -- (s2);
    \draw[arr] (s2) -- (s3);
    \draw[arr] (s3) -- (s4);
    \draw[arr] (s4) -- (s5);
    % Transition gauche -> droite
    \draw[arr] (s5.east) -- ++(0.2,0) |- (s6.west);
    % Flèches colonne droite (remonte)
    \draw[arr] (s6) -- (s7);
    \draw[arr] (s7) -- (s8);
    \draw[arr] (s8) -- (s9);
    \draw[arr] (s9) -- (s10);
  \end{tikzpicture}%
  }
  \caption{Cycle de réalisation d'un lot applicatif au sein de VESI.}
  \label{fig:workflow}
\end{figure}

\begin{enumerate}[leftmargin=*, nosep]
  \item \textbf{Spécifications fonctionnelles}~: l'analyste fonctionnel rédige
    un document de spécifications détaillé pour chaque lot, décrivant les règles
    métier, les maquettes d'écran et les critères d'acceptation.
  \item \textbf{Réunion de cadrage}~: l'analyste fonctionnel et le chef de
    projet présentent le document à l'équipe de développement lors d'une réunion
    dédiée, au cours de laquelle les questions techniques sont levées.
  \item \textbf{Développement et conception}~: le développeur implémente la
    solution sur son environnement local (machine personnelle), en respectant
    l'architecture et les conventions du projet.
  \item \textbf{Déploiement en environnement DEV}~: le code est déployé sur les
    ressources Azure de l'environnement de développement (App~Service, base de
    données, Key~Vault, etc.).
  \item \textbf{Tests développeur}~: le développeur vérifie le bon
    fonctionnement dans l'environnement DEV, en testant les scénarios décrits
    dans les spécifications.
  \item \textbf{Mise en recette}~: une fois les tests développeur validés, le
    code est déployé en environnement de recette. L'analyste fonctionnel réalise
    ses propres tests et constitue un répertoire de retours (captures d'écran,
    vidéos) documentant les anomalies et écarts constatés.
  \item \textbf{Correction des retours}~: le développeur corrige les anomalies
    signalées et effectue un nouveau cycle de tests.
  \item \textbf{Tests finaux et validation}~: l'analyste fonctionnel valide les
    corrections en recette. Cette étape conditionne le passage en production.
  \item \textbf{Mise en production}~: le code validé est déployé sur
    l'environnement de production, accessible aux utilisateurs finaux.
  \item \textbf{Monitoring}~: l'équipe surveille l'application en production
    (journaux, alertes, retours utilisateurs) pour détecter d'éventuelles
    régressions.
\end{enumerate}

Ce processus a été appliqué à chacun des lots de \csbag ainsi qu'aux autres
projets présentés ci-après.

% =================================================================
\section{\csbag v2 : conception et réalisation}
\label{sec:csbag}

\subsection{L'existant : \csbag v1 et ses limites}
\label{sec:existant}
% >>> À RÉDIGER <<<

\subsection{Expression des besoins et exigences}
\label{sec:besoins}
% >>> À RÉDIGER <<<

\subsection{Architecture cible}
\label{sec:architecture}
% >>> À RÉDIGER <<<

\subsection{Modèle de données et migration}
\label{sec:donnees-migration}
% >>> À RÉDIGER <<<

\subsection{Réalisation}
\label{sec:realisation}
% >>> À RÉDIGER <<<

\subsection{Qualité, tests et déploiement}
\label{sec:qualite}
% >>> À RÉDIGER <<<

\subsection{Résultats}
\label{sec:resultats-csbag}
% >>> À RÉDIGER <<<

% =================================================================
\section{Elium Sauvegarde}
\label{sec:elium}

Le service Collaborative Solutions utilise la plateforme Elium comme base de
connaissances interne. Afin de prévenir toute perte de contenu et de centraliser
l'information dans l'écosystème SharePoint, un outil de synchronisation
automatique a été développé.

\subsection{Objectif et besoin}
\label{sec:elium-objectif}

L'outil \textit{Elium Sauvegarde} est une application console .NET~6 dont le
rôle est de répliquer les articles (\textit{stories}) publiés sur Elium vers des
pages SharePoint~Online. Il supporte deux modes de synchronisation~: un mode
complet (toutes les stories) et un mode incrémental (basé sur une fenêtre de
jours configurable), permettant une exécution planifiée régulière.

\subsection{Fonctionnement technique}
\label{sec:elium-technique}

Le processus de synchronisation se déroule en plusieurs étapes~:
\begin{enumerate}[nosep]
  \item lecture des éléments de référence existants sur le site SharePoint cible~;
  \item récupération des stories depuis l'API Elium~;
  \item calcul des différences (stories supprimées, modifiées, nouvelles)~;
  \item suppression des pages SharePoint obsolètes et de leurs entrées en liste~;
  \item création et publication des nouvelles pages~;
  \item journalisation détaillée des résultats (succès et erreurs).
\end{enumerate}

L'application s'appuie sur Azure~Key~Vault pour la gestion des secrets
(jetons d'API, certificats) et utilise Microsoft~Graph pour l'interaction avec
SharePoint.

\subsection{Résultats}
\label{sec:elium-resultats}
% >>> À RÉDIGER <<<

% =================================================================
\section{Rapports d'accès TTC/GO}
\label{sec:ttcgo}

La gouvernance des accès aux ressources partagées (sites SharePoint, partages
de fichiers) est un enjeu majeur pour un groupe décentralisé comme VINCI~Energies.
Le projet \textit{Access~Report~TTC/GO} automatise la génération de ces rapports
de permissions.

\subsection{Objectif et besoin}
\label{sec:ttcgo-objectif}

Ce système repose sur plusieurs Azure~WebJobs qui traitent des demandes de
rapports d'accès. Un job de gouvernance alimente des files d'attente
(Azure~Storage~Queue) avec les demandes entrantes, puis des jobs spécialisés
(TTC, GO, fichiers BU, confidentiels) exécutent les exports de permissions
correspondants.

\subsection{Fonctionnement technique}
\label{sec:ttcgo-technique}

Le traitement d'une demande suit le parcours suivant~:
\begin{enumerate}[nosep]
  \item réception d'une requête depuis la file de gouvernance~;
  \item vérification des doublons et des demandes déjà traitées~;
  \item export des permissions (NTFS ou SharePoint selon le périmètre)~;
  \item génération d'un rapport Excel~;
  \item envoi du rapport par courriel (en pièce jointe ou via un lien de
    partage pour les fichiers volumineux)~;
  \item journalisation de l'exécution en base SQL~Server.
\end{enumerate}

L'application utilise le SDK Azure~WebJobs avec des déclencheurs sur file
d'attente et sur minuteur, Azure~Key~Vault pour les secrets, et
Application~Insights pour le monitoring.

\subsection{Résultats}
\label{sec:ttcgo-resultats}
% >>> À RÉDIGER <<<
